\section{연구 방법 및 절차}

\subsection{모델의 범위와 가정}
본 연구는 \texttt{model\_me.m}의 다리 운동 모형과 \texttt{optimize\_chainring.m}의 수치 최적화를 사용한다. 체인링이 바꾸는 고관절 각속도와 요구 토크를 계산하고, 그 순간 낼 수 있는 최대 토크에 대한 사용률을 부담으로 정의한다. 주요 목표는 같은 주행조건에서 부담의 변동을 줄이는 것이다.

체인링과 다리 마디 등 모든 대상은 강체이며 체인은 질량이 없고 마찰이 무한대인 이상적인 줄이라 가정했다. 이후의 가정은 내용이 나오면 하는 것이 좋을 것 같다.
\begin{enumerate}
\item 체인링과 다리 마디는 강체이고 체인은 늘어나거나 미끄러지지 않는다. 체인 방향은 일정하며, 체인링의 형상은 볼록하다.
\item 고관절 A와 크랭크축 O는 고정한다. 다리는 균일한 막대 두 개로 나타내고, 능동 토크는 고관절에서만 발생한다. 다리의 중력에 의한 위치에너지 변화와 관절 손실은 무시한다.
\item 한 시점에는 페달을 미는 다리 하나만 계산한다. 반 바퀴마다 활성 다리를 교대하여 다리의 기하와 관성을 \(\pi\) 주기로 확장한다. 반대쪽 다리의 양의 구동 기여와 수동 저항은 이 단순화에서 계산하지 않는다.
\item 주행속도 \(v_0>0\)와 저항력 \(f\geq0\)는 고정하고, 뒷바퀴와 스프로킷은 함께 회전한다고 가정한다. 따라서 체인 속도 \(V_0=r v_0/R\)는 일정하다.
\item \(k=d\phi/d\theta\)가 0에 가까운 짧은 구간은 관성으로 통과하고 속도 변화는 무시한다고 근사한다. 그 구간은 부담 평가에서 제외한다. 이 가정은 별도의 동역학 검증을 완료했다는 뜻은 아니다.
\end{enumerate}


\begin{table}[!htb]
\centering\small
\begin{tabular}{@{}lll@{}}
\toprule
기호 & 의미 & 예시 값\\\midrule
\(L_1,L_2\) & 허벅지, 정강이 길이 & \(0.44,\ 0.45\) m\\
\(\ell,d\) & 크랭크 길이, 고관절 높이 & \(0.17,\ 0.62\) m\\
\(m_1,m_2\) & 허벅지, 정강이와 발의 질량 & \(8,\ 4\) kg\\
\(I_c\) & 크랭크 및 체인링 관성 & \(0.05\) kg\(\cdot\)m\(^2\)\\
\(\bar\rho,r,R\) & 평균 유효 반지름, 스프로킷 및 바퀴 반지름 & \(0.083,\ 0.025,\ 0.34\) m\\
\(\tau_{\mathrm{iso}},\omega_{\max}\) & 최대 토크 모형의 매개변수 & \(200\) N\(\cdot\)m, \(15\) rad/s\\
\(N,N_H\) & 각도 격자 수, 푸리에 최대 차수 & \(1440,\ 3\)\\
\(k_{\min}\) & 관성 통과 구간의 기준 & \(0.05\)\\
\bottomrule
\end{tabular}
\caption{계산에 사용할 매개변수의 예시. 인체 능력과 관성의 수치는 측정값이 아닌 가정값이다.}
\label{tab:param}
\end{table}
\texttt{model\_me.m}의 \(v_0,f,r,R\) 초기값 0은 미설정 표시이다. 최적화 함수는 기본값을 자동으로 넣지 않으므로 계산 전에 필요한 값을 직접 정한다.

\subsection{체인링과 모멘트팔}
\begin{figure}[!tb]
\centering\includegraphics[width=0.76\textwidth]{f_geom.pdf}
\caption{모델의 기하 (\(\theta=290^\circ\), 실제 비율). (a) 고관절 A, 무릎 B, 페달 C, 크랭크축 O와 각 \(\alpha,\beta,\theta\). 그림에서 자전거 앞은 왼쪽이다. (b) 모멘트팔 \(\rho_\pm\)는 축에서 체인 줄까지 내린 수직선의 길이이다.}
\label{fig:geom}
\end{figure}
크랭크축 O를 원점, 자전거 뒤쪽을 \(+x\)로 두고, 크랭크 각 \(\theta\)는 크랭크가 맨 아래일 때 0, 앞으로 달리는 방향이 \(+\)이다([Fig.~\ref{fig:geom}]). 체인링 모양은 체인 핀 중심이 지나는 피치곡선으로 나타낸다. 토크의 팔은 축에서 힘의 작용선까지의 수직거리이므로, 체인 장력의 팔은 축에서 접점까지의 거리 \(R\)이 아니라 축에서 체인 줄까지의 수직거리이다. 위쪽, 아래쪽 체인 줄까지의 수직거리를 모멘트팔 \(\rho_+\), \(\rho_-\)라 한다. 원형에서는 \(\rho_\pm=R\)이지만 비원형에서는 다르다.



% ===== [D2] 모멘트팔 계산: 시작 =====
\begin{dupA}{D2}{접점의 방향과 수직거리}
체인링 몸체에 붙은 극좌표에서 피치곡선을 \(R(\varphi)\)로 쓴다. 접점에서 반지름 방향과 곡선의 법선이 이루는 각을 \(\gamma\)라 하면, 축에서 접선까지의 수직거리는 반지름을 법선 방향으로 내린 길이이다.
\[
\rho=R\cos\gamma .
\]
원이면 \(\gamma=0\)이라 \(\rho=R\)이고, 반지름이 각도에 따라 빨리 변할수록 \(\rho\)는 \(R\)보다 작아진다.
\end{dupA}
\begin{dupB}{D2}{극좌표 형상에서 모멘트팔}
이론 문서는 접선벡터 \(\frac{d}{d\varphi}(R\mathbf e_r)=\frac{dR}{d\varphi}\mathbf e_r+R\,\mathbf e_\varphi\)와 위치벡터의 외적으로 모멘트팔을 구하였다.
\[
\rho=\left|R\mathbf e_r\times\frac{(dR/d\varphi)\mathbf e_r+R\mathbf e_\varphi}{\sqrt{(dR/d\varphi)^2+R^2}}\right|=\frac{R^2}{\sqrt{R^2+(dR/d\varphi)^2}} .
\]
접선벡터에서 반지름 방향 성분이 \(dR/d\varphi\), 수직 성분이 \(R\)이므로 \(\tan\gamma=(dR/d\varphi)/R\), 즉 \(\cos\gamma=R/\sqrt{R^2+(dR/d\varphi)^2}\)이다. 따라서 이 식은 가의 \(\rho=R\cos\gamma\)와 같다.
\end{dupB}
% ===== [D2] 끝 =====

\subsection{속도의 연결과 운동방정식}
% ===== [D1] 속도 연결: 시작 =====
\begin{dupA}{D1}{같은 체인은 같은 속도}
위쪽 체인은 한 줄의 곧은 줄이므로 앞뒤 끝에서 줄 방향 속도가 같다. 뒤 스프로킷은 원이라 체인 속도가 \(r_s\omega_s\)이다. 앞에서는 팔이 \(\rho_+\)인 지레가 돌리는 것과 같으므로 체인 속도가 \(\rho_+\omega\)이다. 따라서 자전거 속도가 일정해도 크랭크 각속도는
\[
\omega(\theta)=\frac{r_s\,\omega_s}{\rho_+(\theta)}
\]
로, 팔이 긴 각도에서 느리고 짧은 각도에서 빠르다.
\end{dupA}
\begin{dupB}{D1}{접선 방향 성분으로 쓴 속도 조건}
접점의 회전 속도는 위치벡터에 수직이고 크기가 \(R\omega\)이다. 이를 체인 방향으로 사영하면 위치벡터에서 접선에 수직인 성분, 즉 \(\rho_+\)에 \(\omega\)를 곱한 값이 된다. 뒤 스프로킷에서는 같은 성분이 \(r_s\omega_s\)이므로 아래 식 (\ref{eq:1})을 얻는다. 가의 \(\omega=r_s\omega_s/\rho_+\)는 이 식을 \(\omega\)에 대해 푼 것이다.
\end{dupB}
% ===== [D1] 끝 =====
\begin{equation}
\rho_+(\theta)\,\omega=r_s\,\omega_s \label{eq:1}
\end{equation}
식 (\ref{eq:1})을 시간에 대해 미분하고 \(d\rho_+/dt=(d\rho_+/d\theta)\,\omega\)를 쓰면
\begin{equation}
\dot\omega_s=\frac{\rho_+}{r_s}\,\dot\omega+\frac{1}{r_s}\frac{d\rho_+}{d\theta}\,\omega^2 . \label{eq:2}
\end{equation}
앞축에는 회전 운동의 제2법칙을, 뒤축에는 바퀴 회전과 자전거 전체의 병진을 합친 일-에너지 정리를 적용한다. \(I_c\)는 크랭크와 체인링의 관성, \(I_r=I_w+MR_w^2\)는 바퀴 관성 \(I_w\)와 사람과 자전거의 질량 \(M\)을 뒤축 회전으로 환산한 관성이다.
\begin{align}
I_c\,\dot\omega&=\tau_p-\rho_+T_++\rho_-T_- , \label{eq:3}\\
I_r\,\dot\omega_s&=r_s\,(T_+-T_-)-\tau_s . \label{eq:4}
\end{align}
식 (\ref{eq:4})에서 장력 차를 구해 식 (\ref{eq:3})에 넣고 식 (\ref{eq:2})를 대입하면 장력이 사라진 운동방정식을 얻는다.
\begin{equation}
\Bigl(I_c+I_r\frac{\rho_+^2}{r_s^2}\Bigr)\dot\omega=\tau_p-\frac{\rho_+}{r_s}\tau_s-I_r\frac{\rho_+}{r_s^2}\frac{d\rho_+}{d\theta}\,\omega^2+T_-(\rho_--\rho_+) \label{eq:5}
\end{equation}
왼쪽에는 뒤쪽 관성이 \((\rho_+/r_s)^2\)배로 더해지고, 오른쪽 셋째 항은 팔이 변할 때 뒤쪽을 가속하는 데 드는 토크이다.

\subsection{다리 모델과 페달 토크}
\begin{figure}[!tb]
\centering\includegraphics[width=0.76\textwidth]{f_leg.pdf}
\caption{오른다리의 힘 (\(60\) rpm, \(\theta=260^\circ\)). (a) 고관절 토크가 허벅지를 돌리면 정강이가 페달을 정강이 방향으로 민다. (b) 페달 힘 \(F\)를 크랭크에 수직인 \(F_t\)와 크랭크 방향 \(F_r\)로 나눈다.}
\label{fig:leg}
\end{figure}
고관절 A와 크랭크축 O는 고정하고, 고관절만 토크 \(M_A\)를 내며 무릎과 발목은 핀으로 둔다. 허벅지 AB(\(L_1\))와 정강이 BC(\(L_2\))가 \(+x\)축과 이루는 각을 \(\alpha,\beta\)라 한다.

% ===== [D3] 페달 토크: 시작 =====
\begin{dupA}{D3}{힘과 팔로 보는 페달 토크}
정강이는 무릎 B와 페달 C 두 점에서만 힘을 받는다. 무릎은 토크를 내지 않으므로 B를 중심으로 정강이의 토크를 더하면 C의 힘이 만드는 토크만 남고, 이것이 0이어야 한다. 따라서 페달에 걸리는 힘은 정강이 방향이다.
\[
\vec F=F\,\hat u,\qquad \hat u=\frac{\vec C-\vec B}{|\vec C-\vec B|}.
\]
A와 O에서 정강이 작용선까지의 수직거리를 \(d_A,d_O\)라 하면 허벅지의 토크 균형과 크랭크의 토크는
\[
M_A=F\,d_A,\qquad \tau_p=F\,d_O=\ell F_t
\quad\Longrightarrow\quad \tau_p=M_A\,\frac{d_O}{d_A}
\]
이다. \(F_t\)는 \(F\) 중 크랭크에 수직인 성분이다([Fig.~\ref{fig:leg}]).
\end{dupA}
\begin{dupB}{D3}{일의 관계로 쓴 페달 토크}
이론 문서는 다리에 에너지가 쌓이지 않으므로 고관절이 한 일과 크랭크로 간 일이 같다는 가상일 원리를 썼다\cite{lynch2017}.
\[
M_A\,d\alpha=\tau_p\,d\theta\quad\Longrightarrow\quad \tau_p=M_A\frac{d\alpha}{d\theta}.
\]
가와 같은지 확인한다. 정강이는 휘지 않으므로 B와 C의 속도를 정강이 방향 \(\hat u\)로 사영한 값이 같다. B는 A를 중심으로 각속도 \(\dot\alpha\)로, C는 O를 중심으로 \(\omega\)로 돌므로 그 사영은 각각 \(\dot\alpha\,d_A\), \(\omega\,d_O\)이다. 따라서 \(\dot\alpha\,d_A=\omega\,d_O\), 즉 \(d\alpha/d\theta=d_O/d_A\)이다. 1440개 각도에서 두 방법의 토크 차이는 \(10^{-14}\) N\(\cdot\)m 이하였다.
\end{dupB}
% ===== [D3] 끝 =====
두 다리에 각각 적용해 더한다. 다리 좌표식을 미분하면
\begin{equation}
\tau_p=M_{A}\frac{d\alpha}{d\theta},\qquad \frac{d\alpha}{d\theta}=\frac{\ell\cos(\beta-\theta)}{L_1\sin(\beta-\alpha)} \label{eq:6}
\end{equation}
\(d\alpha/d\theta=0\)인 데드스팟은 오른다리 \(\theta=39.0^\circ\), \(221.5^\circ\)이며, 이때 정강이 작용선이 크랭크축을 지난다.

\subsection{체인링 형상과 크랭크 운동}
누적 크랭크 각도를 \(\theta(t)\)라 하고, 체인 접선까지의 수직거리인 유효 반지름을 \(\rho(\theta)\)라 한다. 코드의 \texttt{rho\_function}과 \texttt{rho\_prime}은 각각 다음 식과 그 해석적 미분을 평가한다.
\begin{align}
\rho(\theta)&=\bar\rho\left[1+\sum_{n=1}^{N_H}\{a_n\cos(n\theta)+b_n\sin(n\theta)\}\right],\label{eq:14}\\
\rho'(\theta)&=\bar\rho\sum_{n=1}^{N_H}n\{-a_n\sin(n\theta)+b_n\cos(n\theta)\}.\label{eq:me-rhop}
\end{align}
\(a_n,b_n\)은 무차원 계수이며, 상수항을 고정하여 평균 유효 반지름을 \(\bar\rho\)로 유지한다.
체인 속도 일치 조건에서
\begin{equation}
\dot\theta=\frac{V_0}{\rho},\qquad
\ddot\theta=-\frac{V_0^2\rho'}{\rho^3},\qquad V_0=\frac{r v_0}{R}
\label{eq:me-motion}
\end{equation}
를 얻는다. \texttt{theta}는 \(\theta(0)=0\)을 초기조건으로 \texttt{ode45}를 사용한다. 최적화는 각도별로 식 (\ref{eq:me-motion})을 직접 평가하므로 후보 형상마다 미분방정식을 다시 풀 필요가 없다. 누적 각도와 체인링 형상은 \(\pi\)로 접지 않는다. 따라서 다리 모형은 \(\pi\) 주기이지만, 일반 푸리에 체인링의 주기는 \(2\pi\)일 수 있다.

\subsection{활성 다리의 기하와 관성}
활성 다리의 각도는 \(\vartheta=\operatorname{mod}(\theta,\pi)\)로 둔다. \texttt{fourbar}에서 다음을 정의한다.
\[
P=2L_1(d+\ell\cos\vartheta),\quad Q=2L_1\ell\sin\vartheta,\quad
K=L_1^2+\ell^2+d^2+2d\ell\cos\vartheta-L_2^2,
\]
\texttt{fourbar}가 사용하는 기하학적 해와 각도비는
\begin{equation}
\phi=\operatorname{atan2}(Q,P)+\arccos\frac{K}{\sqrt{P^2+Q^2}},\qquad
k=\frac{K'-P'\cos\phi-Q'\sin\phi}{Q\cos\phi-P\sin\phi}
\label{eq:me-fourbar}
\end{equation}
이다. 여기서 프라임은 활성 다리 각도에 대한 미분이다. \(P'=-2L_1\ell\sin\vartheta\), \(Q'=2L_1\ell\cos\vartheta\), \(K'=-2d\ell\sin\vartheta\)를 사용한다.

다리의 운동에너지를 \(T_{\mathrm{leg}}=\tfrac12J(\vartheta)\dot\theta^2\)로 쓰면 \texttt{legJ}의 등가 관성은
\begin{equation}
J=\frac{m_1L_1^2}{3}k^2+
\frac{m_2}{4}\{L_1^2k^2+\ell^2+2L_1\ell k\cos(\phi-\vartheta)\}
+\frac{m_2L_2^2}{12}\lambda_{\mathrm{leg}}^2
\label{eq:me-inertia}
\end{equation}
이다. 정강이의 방향각 변화율 \(\lambda_{\mathrm{leg}}=(uv'-vu')/L_2^2\)에서
\[
\begin{aligned}
u&=\ell\sin\vartheta-L_1\sin\phi,&v&=-\ell\cos\vartheta-d+L_1\cos\phi,\\
u'&=\ell\cos\vartheta-L_1k\cos\phi,&v'&=\ell\sin\vartheta-L_1k\sin\phi
\end{aligned}
\]
를 사용한다. 둘째 항의 \(1/4\)는 정강이 질량중심이 \((B+C)/2\)이기 때문에 생긴다. \(I_c\)는 \(J\)에 포함하지 않고 페달 토크에 한 번만 반영한다.

\(J'\)는 \(h=2\pi/N\)으로 차분한다. 교대 경계에서 반대쪽 다리의 관성과 섞이지 않도록 \(\vartheta_+=\min(\vartheta+h,\pi-\epsilon)\), \(\vartheta_-=\max(\vartheta-h,0)\)으로 제한하여
\[
J'\approx\frac{J(\vartheta_+)-J(\vartheta_-)}{\vartheta_+-\vartheta_-}
\]
로 계산한다. 내부에서는 중앙차분이고 경계 근처에서는 비대칭 차분이다. 교대점에서는 주기 확장한 다리의 관성이 불연속일 수 있으므로 미분을 정의하지 않는다.

\subsection{요구 토크와 순간 부담}
\texttt{tau\_p}는 일정한 주행속도와 저항력에서 필요한 페달 토크를 계산한다.
\begin{equation}
\tau_p=\frac{\rho}{r}fR-I_c\frac{r^2v_0^2}{R^2\rho^3}\rho'.
\label{eq:me-taup}
\end{equation}
다리의 입력 동력과 출력 동력 차이가 운동에너지 변화율과 같다는 관계에서
\begin{equation}
\tau_h=\frac{\tau_p+J\ddot\theta+\tfrac12J'\dot\theta^2}{k}
\label{eq:me-tauh}
\end{equation}
를 얻는다. 이것이 \texttt{leg}의 계산식이다. 토크를 최대 능력에 비례하도록 자유롭게 배분하는 것이 아니라, 형상과 주행조건을 정하면 요구 토크가 결정된다.

최대 고관절 토크를 각속도에 따라 감소하는 직선으로 가정하고 순간 부담을 다음처럼 정의한다.
\begin{equation}
\tau_{\max}(\dot\phi)=\max\left[\tau_{\mathrm{iso}}\left(1-\frac{|\dot\phi|}{\omega_{\max}}\right),0\right],\qquad
u=\frac{|\tau_h|}{\tau_{\max}(\dot\phi)}.
\label{eq:8}
\end{equation}
이는 토크 능력의 사용률이며 실제 피로나 대사 에너지를 직접 측정한 값은 아니다. 토크 방향과 자세에 따른 능력 차이는 생략한다. 최대 토크가 0인데 요구 토크가 0이 아니면 \(u=\infty\), 둘 다 0이면 \(u=0\)으로 정의한다.

\subsection{관성 통과 구간과 목적함수}
\(k=0\)에서는 유한한 고관절 토크가 순간적으로 크랭크에 일을 전달할 수 없다. 본 방법은 \(k_{\min}\)을 최적화 전에 고정하고
\[
\Omega_{\mathrm{coast}}=\{\theta:|k|<k_{\min}\},\qquad
\Omega_{\mathrm{drive}}=\{\theta:|k|\geq k_{\min}\}
\]
로 나눈다. 부담은 구동 구간에서만 평가한다. 교대점도 별도로 제외하며, 관성 통과의 실제 감속과 교대 순간의 에너지 전달은 이 최적화로 검증되지 않는다.

구동 구간의 시간 평균을 \(\langle g\rangle_t\)로 쓰면 \(dt=\rho\,d\theta/V_0\)이므로
\begin{equation}
\langle g\rangle_t=\frac{\int_{\Omega_{\mathrm{drive}}}\rho(\theta)g(\theta)\,d\theta}
{\int_{\Omega_{\mathrm{drive}}}\rho(\theta)\,d\theta},\qquad
\mu_u=\langle u\rangle_t,\quad \sigma_u=\sqrt{\langle(u-\mu_u)^2\rangle_t}.
\label{eq:me-average}
\end{equation}
각도 평균이 아니라 이 시간 평균을 사용한다.

\begin{table}[!htb]
\centering\small
\begin{tabular}{@{}L{19mm} L{41mm} L{80mm}@{}}
\toprule
선택 이름 & 목적함수 & 최소화하는 내용\\\midrule
\texttt{peak} & \(\max u\) & 가장 부담이 큰 순간을 완화한다.\\
\texttt{mean} & \(\mu_u\) & 평균적인 능력 사용률을 낮춘다.\\
\texttt{rms} & \(\sqrt{\langle u^2\rangle_t}\) & 평균 부담과 높은 부담을 함께 줄인다.\\
\texttt{flat} & \(\sigma_u\) & 부담의 변동을 직접 줄인다.\\
\texttt{mixed} & \(\mu_u+\lambda\sigma_u\) & 평균과 평탄도의 비중을 \(\lambda\geq0\)으로 조절한다.\\
\bottomrule
\end{tabular}
\caption{비교할 수 있는 부담 목적함수 후보. 모든 최대값과 평균은 구동 구간에서 계산한다.}
\label{tab:me-objectives}
\end{table}

현재 코드는 \texttt{flat}만 사용한다. 부담이 높지만 일정한 해를 막기 위해 \(\mu_u\leq\mu_{u,\mathrm{circle}}\)을 함께 둔다. 이는 추가 설계 조건이며 같은 일을 한다는 사실만으로 자동 성립하지 않는다. 후보 \texttt{rms}는
\[
\langle u^2\rangle_t=\mu_u^2+\sigma_u^2
\]
를 줄이지만 표준편차만의 최소화를 보장하지는 않는다.

현재 코드는 최대 부담이 1 이하라는 조건을 강제하지 않는다.

\subsection{푸리에 계수의 수치 최적화}
\subsubsection{형상 제약}
\(\rho\)는 공간에서 고정된 체인 방향에 대한 지지함수를 회전각으로 나타낸 값이다. 따라서 볼록 피치곡선을 만들려면
\begin{equation}
\rho+\rho''>0
\label{eq:13}
\end{equation}
이어야 한다. 기본 설정에서는
\[
0.45\bar\rho\leq\rho\leq1.55\bar\rho,\qquad
\rho+\rho''\geq0.4\bar\rho
\]
를 사용한다. 이 조건은 지지함수의 최댓값/최솟값 비율을 제한하며, 임의 형상에서 실제 반지름의 지름비와 같은 뜻으로 사용하지 않는다.

\subsubsection{이산화와 풀이}
한 바퀴를 짝수 \(N\)개의 중점 \(\theta_j=2\pi(j+1/2)/N\)으로 나누어 교대점 \(0,\pi\)를 피한다. 기하 \(k_j\), 관성 \(J_j\), 미분 \(J'_j\)와 구동 구간 표시는 형상에 무관하므로 한 번만 계산한다. 각 계수 후보에서는 다음 순서로 평가한다.
\begin{enumerate}
\item 푸리에 식으로 \(\rho_j,\rho'_j\)를 계산한다.
\item 식 (\ref{eq:me-motion})으로 \(\dot\theta_j,\ddot\theta_j\)를 구한다.
\item 식 (\ref{eq:me-taup}), (\ref{eq:me-tauh})으로 \(\tau_{p,j},\tau_{h,j}\)를 계산한다.
\item 구동 구간에서 최대 토크와 부담 \(u_j\)를 구한다.
\item 시간 가중치 \(w_j=\rho_j/\sum_{i\in\Omega_{\mathrm{drive}}}\rho_i\)로 평균과 표준편차를 계산한다.
\end{enumerate}
기본 최적화 문제는
\begin{equation}
\begin{aligned}
\min_{a_1,b_1,\ldots,a_{N_H},b_{N_H}}\quad &\sigma_u\\
\text{s.t.}\quad &\mu_u\leq\mu_{u,\mathrm{circle}},\\
&\rho_{\min}\leq\rho(\theta)\leq\rho_{\max},\quad
\rho(\theta)+\rho''(\theta)\geq\kappa_{\min}\bar\rho
\end{aligned}
\label{eq:opt}
\end{equation}
이다. 수치적으로는 구동 격자에서 \(\tau_{\max}/\tau_{\mathrm{iso}}\geq10^{-6}\)을 추가하여 0으로 나누는 일을 방지한다. 탐색 중 분모를 같은 작은 값으로 제한하지만, 이 제약을 만족하는 해에서는 제한 전의 부담 정의와 같다.

형상 제약은 격자 사이도 고려한다. \(A_n=\sqrt{a_n^2+b_n^2}\)일 때 반 격자 간격에서 생길 수 있는 변화의 상한은
\[
m_\rho=\frac{\pi}{N}\sum_n nA_n,\qquad
m_\kappa=\frac{\pi}{N}\sum_n n|1-n^2|A_n
\]
이다. 정규화된 격자값에 이만큼의 여유를 두어 반지름과 곡률 제약을 검사한다. 반면 부담의 최대값과 능력 제약은 격자에서 평가하므로 격자 사이까지 보장하지 않는다.

요구 토크와 최대 토크, 시간 가중치가 형상에 따라 변하므로 이 문제는 볼록 최적화로 취급하지 않는다. MATLAB의 \texttt{fmincon}을 사용하며 원형 계수에서 시작한다. 지역 탐색이므로 전역 최적해를 얻었다고 주장하지 않는다.

\subsection{비교 지표와 검증 범위}
코드는 원형과 최적 형상의 부담 표준편차와 푸리에 계수를 반환한다. 목적값 감소만으로 제외 구간의 감속이나 실제 피로 감소를 검증한 것으로 보지 않는다.

현재 코드에는 체인 장력, 유한 축간거리, 프리휠 운동방정식, 실험 검증이 구현되어 있지 않으므로 이전 코드의 검증 수치나 개선율을 현재 방법의 결과로 옮기지 않는다.

후속 계산에서는 \(N,N_H,k_{\min},\tau_{\mathrm{iso}},\omega_{\max}\)를 바꾸어 계수와 부담 지표의 민감도를 확인한다. 특히 \(k_{\min}\)은 계산상의 작은 오차가 아니라 구동을 생략하는 범위를 정하는 가정이다. 새 방법으로 재계산하기 전까지 기존 결과 절의 수치와 그림은 구형 모형의 결과로 구분한다.
